% !TEX program = pdflatex
% =============================================================================
%  MCM/ICM（美国大学生数学建模竞赛，美赛）论文 LaTeX 模板
% -----------------------------------------------------------------------------
%  编译方式（本模板同时兼容 pdfLaTeX 与 XeLaTeX，二者选一即可）：
%      pdflatex main.tex
%      pdflatex main.tex        % 跑两遍以解析 \ref / \pageref / 目录
%  或：  xelatex  main.tex   （跑两遍）
%  或：  latexmk -pdf main.tex
%
%  !! 重要声明 !!
%  本模板注释中给出的页数上限、页眉格式、身份信息禁令等口径，依据的是
%  COMAP 官方竞赛规则页（MCM/ICM Instructions，现行版本，2027 赛季）：
%      https://www.contest.comap.com/undergraduate/contests/mcm/instructions.html
%  该规则每年都可能修订（例如"代码计入 25 页"就是 2027 赛季的新规），
%  正式参赛前必须以**当年官方发布的通知与 Instructions 页面**为准。
% =============================================================================
%
%  ------------------------- 美赛官方硬性规范速查（2027 赛季口径） -------------------------
%
%   [页数]  **25 页上限覆盖整个提交文件**，包括 Summary Sheet、Solution、
%           Reference List、Table of Contents、Notes、Appendices、Code，
%           以及题目特定的要求（memo / letter / article 等）。
%           ★ 2027 赛季新规：程序代码也算页数！这与国赛"附录页数不限"完全相反。
%   [格式]  必须是 Adobe PDF、英文撰写、字体不小于 12 磅。
%   [页眉]  **每一页的顶部**必须包含队伍控制号和页码，
%           官方示例格式：  Team # 0000000, Page 6 of 25
%   [匿名]  **学生、指导教师、学校名称不得出现在任何一页上**；
%           除控制号外不得含任何身份信息。
%   [首页]  第 1 页是 Summary Sheet（摘要页）。
%   [目录]  Table of Contents 是**被鼓励的**，且**计入** 25 页上限。
%           （注意：这一点与国赛"不要目录"正相反。）
%   [AI]    Report on Use of AI **附在 25 页正文之后**，该部分**无页数限制、
%           不计入 25 页**。AI 使用必须透明：正文内联引用 + 参考文献列出
%           所有用到的 AI 工具。
%   [提交]  PDF 附件须小于 25MB，文件名用控制号（例如 0000000.pdf）。
%           **不要把程序、数据等额外文件随论文提交**。
%
%  =============================================================================
%   ★★★ 最容易无意中超页的地方 ★★★
%   美赛与国赛在附录上的规则**完全相反**：
%       国赛：附录页数不限，鼓励放完整源代码。
%       美赛：附录、代码、目录、参考文献**全部计入 25 页**。
%   因此**不要**在美赛模板里堆砌长篇代码清单或大段数据表。
%   如果确实需要附代码，请把它压缩到最少，或按当年 Instructions 的说明处理。
%   正文每加一节之前，先编译一次并数一遍总页数。
%  =============================================================================

\documentclass[12pt]{article}
% ★ 字体不小于 12 磅 —— 这里是官方硬性要求，请勿改成 11pt/10pt。

% ------------------------------ 数学宏包（★ 必须放在字体之前） ------------------------------
% ★★ 踩坑记录（请不要调整这个顺序）★★
%   amssymb 与 newtxmath 都会定义 \Bbbk、\openbox 等命令。
%   如果先加载 newtxmath 再加载 amssymb，pdfLaTeX 会直接报错：
%       ! LaTeX Error: Command `\Bbbk' already defined.
%       ! LaTeX Error: Command \openbox already defined.
%   正确做法：**先 amsmath/amssymb/amsthm，后 newtxtext/newtxmath**。
\usepackage{amsmath,amssymb,amsthm}
\usepackage{bm}                       % 粗体数学符号 \bm{}

% ------------------------------ 字体（Times 系，12pt） ------------------------------
% newtxtext + newtxmath 是 Times 的现代替代品，字形与数学字体搭配良好。
\usepackage{newtxtext}
\usepackage{newtxmath}
% 备选方案（若 newtx 装不上，任选其一，取消注释即可）：
% \usepackage{times}      % 仅正文 Times，数学字体不变
% \usepackage{mathptmx}   % 正文 + 数学均为 Times（较老，但兼容性最好）

% ---------------------------------- 页面与页眉页脚 ----------------------------------
\usepackage[letterpaper,margin=1in]{geometry}
% 说明：COMAP 未强制纸张尺寸。美式惯例用 letterpaper；
%       若你的排版环境默认 A4，可改成 [a4paper,margin=1in]，不影响合规性。

\usepackage{fancyhdr}
\usepackage{lastpage}

% ★ 官方要求：每一页顶部必须有 "Team # 控制号, Page X of 25"。
%   把下面的 0000000 替换成你的真实队伍控制号（7 位数字，COMAP 分配）。
\newcommand{\MCMTeamID}{0000000}

% ★ 关于 "of 25" 的取舍（重要，请读）：
%   官方示例格式是 "Team # 0000000, Page 6 of 25"，即分母固定写 25。
%   本模板**默认写死为 25**，理由：
%     (a) 与官方示例完全一致，评阅人看到的格式不会意外；
%     (b) 25 是规则上限，写实际页数（如 18）反而可能被误读为"未达要求"。
%   如果你更希望显示实际总页数，把下面一行改为 \MCMUseActualPageCounttrue，
%   它会用 lastpage 的 \pageref{LastPage} 自动取总页数（需要编译两遍）。
%   ★ 无论用哪种写法，**编译后都必须自己数一遍实际页数**，确认没超过 25 页：
%       看编译日志里的 "Output written on main.pdf (N pages)."
%     或命令行执行  pdfinfo main.pdf  （查看 Pages 字段）
\newif\ifMCMUseActualPageCount
\MCMUseActualPageCountfalse          % ← 改成 \MCMUseActualPageCounttrue 即显示实际页数

\ifMCMUseActualPageCount
  \newcommand{\MCMTotalPages}{\pageref{LastPage}}
\else
  \newcommand{\MCMTotalPages}{25}
\fi

\newcommand{\MCMHeaderText}{Team \# \MCMTeamID, Page \thepage\ of \MCMTotalPages}

\pagestyle{fancy}
\fancyhf{}
\fancyhead[C]{\MCMHeaderText}        % 居中：Team # 0000000, Page 6 of 25
\renewcommand{\headrulewidth}{0pt}   % 不画页眉横线（官方未要求，画了也不违规）
\renewcommand{\footrulewidth}{0pt}
% 想让控制号靠左、页码靠右？改用下面两行，并注释掉上面的 \fancyhead[C]：
% \fancyhead[L]{Team \# \MCMTeamID}
% \fancyhead[R]{Page \thepage\ of \MCMTotalPages}

% \tableofcontents 等环境会强制切到 plain 样式，这里重定义 plain 以保持页眉一致
\fancypagestyle{plain}{%
  \fancyhf{}%
  \fancyhead[C]{\MCMHeaderText}%
  \renewcommand{\headrulewidth}{0pt}%
  \renewcommand{\footrulewidth}{0pt}%
}

% ---------------------------------- 表格 ----------------------------------
\usepackage{booktabs}                 % ★ 三线表
\usepackage{array}
\usepackage{multirow}
\usepackage{siunitx}
\sisetup{detect-all=true}

% ---------------------------------- 图形与浮动体 ---------------------------------
\usepackage{graphicx}
\usepackage{float}
\usepackage{caption}
\captionsetup{
  font          = small,
  labelsep      = colon,              % 英文习惯："Figure 1: ..."
  justification = centering,
  skip          = 6pt,
}
% 图题在图下方、表题在表上方，靠 \caption 的书写位置控制。
% \caption 必须写在 \label 之前，否则 \ref 会取到上一个编号。
% ★ 同一结果不要同时用表格和图片呈现 —— 二者重复会浪费宝贵的 25 页配额。

% ---------------------------------- 列表与版式 ----------------------------------
\usepackage{enumitem}
\setlist{topsep=2pt,itemsep=1pt,parsep=0pt}
\usepackage{setspace}
% 美赛不要求特定行距；默认单倍行距最省页数。
% 若指导老师要求双倍行距，取消下一行注释（但请留意页数会大幅增加）：
% \doublespacing

\usepackage{indentfirst}

% ---------------------------------- 代码清单（可选） ---------------------------------
% ★ 警告：在美赛里，附录中的代码**计入 25 页上限**（2027 赛季起明确）。
%   本模板只保留一个极短的示例，用来演示怎么排版；
%   正式论文中请谨慎添加，能不放就不放。
\usepackage{listings}
\usepackage{xcolor}
\lstset{
  basicstyle       = \ttfamily\footnotesize,
  keywordstyle     = \color[RGB]{0,0,180}\bfseries,
  commentstyle     = \color[RGB]{0,128,0}\itshape,
  stringstyle      = \color[RGB]{170,0,0},
  numbers          = left,
  numberstyle      = \tiny\color{gray},
  numbersep        = 8pt,
  frame            = single,
  breaklines       = true,
  showstringspaces = false,
  tabsize          = 4,
  columns          = flexible,
}

% ---------------------------------- 超链接 ----------------------------------
\usepackage[hidelinks]{hyperref}
\hypersetup{
  bookmarksnumbered = false,
  pdftitle  = {MCM/ICM Solution Paper},
  pdfauthor = {Team \MCMTeamID},      % ★ 只能写控制号，不得写姓名/学校
}
\renewcommand{\contentsname}{Table of Contents}

% =============================================================================
\begin{document}

% =============================================================================
% PAGE 1 — SUMMARY SHEET（官方规定第 1 页就是摘要页，必须单独一页）
% -----------------------------------------------------------------------------
%  ★ 本页是评阅人最先、也可能唯一细读的一页，务必把"问题是什么、你用了什么方法、
%    得到什么结果、结论有什么价值"讲清楚，不要只写背景。
%  ★ 官方第六条精神：本页不得出现学生姓名、指导教师姓名、学校名称，
%    只能出现队伍控制号（已由页眉统一给出一遍，正文里不必重复）。
%  ★ 摘要页是否计入 25 页？计入 —— 它是"整个提交文件"的一部分。
%  ★ 同样地：不要在这里放致谢、不要写学校名、不要放校徽。
% =============================================================================

\begin{center}
  {\large\bfseries
   Here Goes the Title of Your MCM/ICM Solution Paper}\\[6pt]
  % ★ 标题替换为你的论文标题。建议具体、可检索，避免"Analysis of a Problem"这类空泛标题。
  {\normalsize\bfseries Summary}
\end{center}

\vspace{4pt}

% ---- Summary 正文（占位） ----
% ★ 建议 350-600 词，占满但不溢出一页。写成 4-6 段，每段一个要点：
%   段 1：问题背景与我们的总体思路
%   段 2：模型一 / 问题一的做法与结果（给出**具体数字**）
%   段 3：模型二 / 问题二的做法与结果
%   段 4：模型三 / 问题三的做法与结果
%   段 5：灵敏度与稳健性检验结论
%   段 6：模型特色与推广价值
We address the problem of \emph{[state the task in one sentence]}. Our approach combines
[method A] with [method B] to produce a quantitative, reproducible decision framework.
Unlike purely descriptive treatments, we couple a mechanistic model with a data-driven
calibration step, and we explicitly quantify uncertainty rather than reporting a single
point estimate. The deliverable is a ranked set of operational recommendations together
with the confidence we attach to each one.

For the first task, we clean the provided data by removing outliers with a $3\sigma$
rule and imputing missing entries by linear interpolation, then screen candidate
predictors using Pearson and Spearman correlation matrices. We compare a multiple linear
regression against a random forest, and select the latter because its ten-fold
cross-validated coefficient of determination reaches $0.93$, against $0.71$ for the
linear model. This model supplies the calibrated inputs used in all later sections.

For the second task, we formulate the decision problem as a mixed-integer nonlinear
program with $180$ binary and $96$ continuous variables, minimizing total cost subject to
capacity, inventory-balance, demand-satisfaction, and logical-exclusivity constraints.
Because the problem is NP-hard, we solve it in two stages: a Lagrangian relaxation
furnishes a lower bound, and an adaptive large neighborhood search drives the upper bound
downward. The algorithm converges after $1200$ iterations with an optimal objective of
$2.417\times10^{6}$ CNY, an $18.6\%$ improvement over a greedy baseline.

For the third task, we allow demand and price to vary within a box uncertainty set and
build a two-stage robust optimization model. Strong duality converts the inner max-min
problem into a tractable linear program. The resulting plan raises expected cost by only
$4.2\%$ under $\pm 15\%$ demand perturbation, and no scenario in our test set violates a
constraint. Because the robust counterpart remains a linear program, it can be solved to
global optimality with a standard solver rather than a heuristic.

We validate the framework on three levels. Residual analysis shows approximately normal
errors with mean $0.03$ and standard deviation $0.87$. One-at-a-time sensitivity analysis
identifies unit transportation cost as the dominant driver, with an elasticity of $1.34$.
A ten-thousand-sample Monte Carlo study bounds the relative variation of the optimal
objective at $6.8\%$ with $95\%$ confidence. The principal strength of our approach is
that it fuses mechanistic constraints with data-driven estimation while representing
uncertainty explicitly, so the recommendations remain defensible when the inputs shift.
On the strength of these results we recommend consolidating shipments on the three
highest-volume arcs and holding one node in reserve; that single change captures most of
the achievable saving while preserving the robustness margin.

\vspace{6pt}
\noindent{\bfseries Keywords:} mixed-integer programming; robust optimization; adaptive
large neighborhood search; sensitivity analysis; Monte Carlo simulation

\newpage

% =============================================================================
% PAGE 2 — TABLE OF CONTENTS
% -----------------------------------------------------------------------------
%  ★ 美赛**鼓励**使用 Table of Contents（与国赛"不要目录"的要求相反）。
%  ★ 目录页**计入** 25 页上限，所以不要写冗长的三级目录；本模板只用两级。
%  ★ 目录需要编译两遍才会正确（第一遍显示 ??）。
% =============================================================================
\tableofcontents
% 只想显示到二级标题（更省页数）：
% \setcounter{tocdepth}{1}

\newpage

% =============================================================================
% 正文 —— 以下各节均计入 25 页
% =============================================================================

\section{Introduction}
\label{sec:intro}

Restate the problem in your own words, then state precisely what you are asked to
deliver. A useful pattern is three short subsections: background, problem statement, and
our approach.

\subsection{Background}

Explain why the problem matters and what is already known. Keep this brief: reviewers
care far more about what you did than about the history of the field.

\subsection{Problem Statement}

Enumerate the required tasks. Keep the numbering identical to the numbering used in
\S\ref{sec:model} onward, so that a reviewer can always map task to method to result.

\subsection{Our Approach}

Give a one-paragraph roadmap of the paper. Mention the modelling philosophy and the
solution strategy in plain language before any mathematics appears.

\section{Assumptions and Justifications}
\label{sec:assumptions}

% ★ 美赛评阅非常看重"假设是否有依据"。每条假设后面都要跟一句理由，
%   并说明该假设在后续哪些模型/公式中起作用。
\begin{enumerate}
  \item \textbf{The supplied data are accurate and free of systematic recording bias.}
        Justification: the data are the only empirical basis available, and no
        independent measurement is provided for cross-checking.
  \item \textbf{External parameters do not undergo structural breaks within the study
        horizon.} Justification: the horizon is short relative to the observed
        historical variability of these parameters.
  \item \textbf{Transportation time is proportional to distance.} Justification:
        routing is over a highway network where average speed varies little between
        the arcs considered.
  \item \textbf{Demand is met in full in every period.} Justification: the problem
        statement treats unmet demand as infeasible rather than as a penalized
        shortfall.
\end{enumerate}

\section{Notations}
\label{sec:notation}

% ★ 表题在表上方；三线表只用 \toprule \midrule \bottomrule，禁止竖线。
% ★ 正文必须用 \ref 引用表格（见表 \ref{tab:notation}），不能只放一张孤表。
Table~\ref{tab:notation} collects the symbols used throughout the paper; symbols not
listed here are defined where they first appear.

\begin{table}[htbp]
  \centering
  \caption{Notation}
  \label{tab:notation}
  \begin{tabular}{cll}
    \toprule
    \textbf{Symbol} & \textbf{Description} & \textbf{Unit / Range} \\
    \midrule
    $i,\ j$          & Index of supply node / demand node          & --- \\
    $t$              & Index of time period                        & --- \\
    $x_{ijt}$        & Quantity shipped from $i$ to $j$ in period $t$ & tonne \\
    $y_{it}$         & Binary: 1 if supply node $i$ is open in $t$ & $\{0,1\}$ \\
    $c_{ij}$         & Unit transportation cost                    & CNY/tonne \\
    $d_{jt}$         & Demand at node $j$ in period $t$            & tonne \\
    $S_i$            & Capacity of supply node $i$                 & tonne \\
    $\alpha$         & Half-width of the uncertainty set           & $0.15$ \\
    $T$              & Number of periods in the horizon            & $12$ \\
    \bottomrule
  \end{tabular}
\end{table}

\section{Model Development}
\label{sec:model}

\subsection{Model I: Prediction}

State the model, then justify each term. Number every equation that you later refer to.
Equation~\eqref{eq:objective} gives the objective of the deterministic core model:
% \eqref 会自动补上括号并生成超链接；\ref 则不带括号。
\begin{equation}
  \min\ Z=\sum_{i=1}^{m}\sum_{j=1}^{n}\sum_{t=1}^{T} c_{ij}\,x_{ijt}
            +\sum_{i=1}^{m}\sum_{t=1}^{T} f_i\,y_{it}
  \label{eq:objective}
\end{equation}
where $f_i$ is the fixed opening cost of supply node $i$. The first sum is variable
transportation cost and the second is fixed opening cost.

The constraints are:
\begin{equation}
  \sum_{i=1}^{m} x_{ijt} = d_{jt}, \qquad \forall j,\ t
  \label{eq:demand}
\end{equation}
\begin{equation}
  \sum_{j=1}^{n} x_{ijt} \le S_i\,y_{it}, \qquad \forall i,\ t
  \label{eq:capacity}
\end{equation}
\begin{equation}
  y_{it}\in\{0,1\},\qquad x_{ijt}\ge 0
  \label{eq:domain}
\end{equation}
Equation~\eqref{eq:demand} enforces demand satisfaction, \eqref{eq:capacity} enforces
capacity, and \eqref{eq:domain} fixes the variable domains.

\subsection{Model II: Optimization}

Present the extended model. For multi-line derivations use \texttt{align}, numbering only
the lines you need:
\begin{align}
  \mathcal{L}(x,y,\lambda,\mu)
    &= \sum_{i,j,t} c_{ij}x_{ijt} + \sum_{i,t} f_i y_{it}
       + \sum_{j,t}\lambda_{jt}\Bigl(d_{jt}-\sum_i x_{ijt}\Bigr) \notag \\
    &\quad + \sum_{i,t}\mu_{it}\Bigl(\sum_j x_{ijt}-S_i y_{it}\Bigr)
  \label{eq:lagrange}
\end{align}

\subsection{Solution Algorithm}

Describe the algorithm, its parameters, its convergence criterion, and the computing
environment. Report the iteration count and the optimality gap.

\section{Results}
\label{sec:results}

% ★ 图题在图下方：先 \includegraphics（或占位框），再 \caption，最后 \label。
% ★ 图与表不要重复呈现同一组数据：表 \ref{tab:results} 已经给出各期成本，
%   就不要再画一张把同一组数字重画一遍的柱状图。
\begin{figure}[htbp]
  \centering
  % ==== 替换为真实图片：\includegraphics[width=0.7\linewidth]{figures/convergence.pdf} ====
  \framebox[0.7\linewidth]{\rule{0pt}{4.5cm}\makebox[0pt][c]{\small Figure placeholder}}
  \caption{Convergence of the adaptive large neighborhood search.}
  \label{fig:convergence}
\end{figure}

As Figure~\ref{fig:convergence} shows, the objective stabilizes after roughly $1200$
iterations with a relative gap below $0.5\%$.

\begin{table}[htbp]
  \centering
  \caption{Optimal cost and demand satisfaction by period}
  \label{tab:results}
  \begin{tabular}{cccc}
    \toprule
    \textbf{Period} & \textbf{Cost (CNY)} & \textbf{Satisfaction (\%)} & \textbf{Open nodes} \\
    \midrule
    1 & $1.83\times10^{5}$ & 100.0 & 3 \\
    2 & $1.96\times10^{5}$ & 100.0 & 3 \\
    3 & $2.11\times10^{5}$ & 100.0 & 4 \\
    \textbf{Total} & $\bm{2.417\times10^{6}}$ & 100.0 & --- \\
    \bottomrule
  \end{tabular}
\end{table}

Table~\ref{tab:results} reports the period-by-period optima. The number of open nodes
rises from three to four as demand grows, consistent with economies of scale.

\section{Sensitivity Analysis}
\label{sec:sensitivity}

% ★ 美赛评阅几乎必然检查灵敏度分析。至少要有一个参数扰动表或图，
%   并指出"哪个参数最敏感、所以现实中应该优先控制它"。
We perturb each key parameter by $\pm 10\%$ one at a time and record the relative change
in the optimal objective.

\begin{table}[htbp]
  \centering
  \caption{Sensitivity of the optimal objective (relative change, \%)}
  \label{tab:sensitivity}
  \begin{tabular}{lccc}
    \toprule
    \textbf{Parameter} & \textbf{$-10\%$} & \textbf{Baseline} & \textbf{$+10\%$} \\
    \midrule
    Unit transportation cost $c_{ij}$ & $-11.7$ & $0.00$ & $+12.9$ \\
    Demand $d_{jt}$                   & $-8.4$  & $0.00$ & $+9.1$  \\
    Capacity $S_i$                    & $+0.6$  & $0.00$ & $-0.3$  \\
    Uncertainty width $\alpha$        & $+2.2$  & $0.00$ & $+3.8$  \\
    \bottomrule
  \end{tabular}
\end{table}

Table~\ref{tab:sensitivity} shows that the objective is most sensitive to unit
transportation cost (elasticity about $1.34$) and essentially insensitive to capacity,
so cost reduction should be prioritized over capacity expansion.

\section{Strengths and Weaknesses}
\label{sec:strengths}

% ★ 这一节是美赛的特色，必须诚实。弱点要具体，并给出改进方向；
%   写"本模型没有缺点"会被视为缺乏批判性思考。
\subsection{Strengths}
\begin{itemize}
  \item The framework couples mechanistic constraints with data-driven calibration,
        so it remains interpretable while capturing nonlinear effects.
  \item The two-stage solution strategy returns high-quality feasible solutions within
        a practical time budget.
  \item Uncertainty is modelled explicitly, so the recommendation does not hinge on a
        single scenario.
\end{itemize}

\subsection{Weaknesses}
\begin{itemize}
  \item Cross-period dependencies such as supplier learning are not modelled, which may
        understate the value of long-term contracts.
  \item The heuristic provides no optimality certificate; the reported gap is empirical.
  \item The data are spatially aggregated, so intra-node heterogeneity is smoothed out.
\end{itemize}

\section{Conclusions}
\label{sec:conclusions}

Summarize what you found, in the same order as the tasks were posed, and give the
concrete numbers. Then state the most important practical implication in one or two
sentences. Do not introduce new material here.

% =============================================================================
% REFERENCES（计入 25 页）
% -----------------------------------------------------------------------------
%  ★ Reference list **计入** 25 页上限。
%  ★ 正文中必须用 \cite{} 引用，且与文献表一一对应。
%  ★ AI 使用必须透明：凡是使用了 AI 工具，都要在正文内联引用，
%    并在本参考文献表中列出所有用到的 AI 工具（见文末 Report on Use of AI）。
% =============================================================================
\newpage
\begin{thebibliography}{99}
  \bibitem{ref1} Giordano F R, Fox W P, Horton S B. A First Course in Mathematical
                 Modeling[M]. 5th ed. Boston: Cengage Learning, 2014.
  \bibitem{ref2} Birge J R, Louveaux F. Introduction to Stochastic Programming[M].
                 2nd ed. New York: Springer, 2011.
  \bibitem{ref3} Ben-Tal A, El Ghaoui L, Nemirovski A. Robust Optimization[M].
                 Princeton: Princeton University Press, 2009.
  \bibitem{ref4} Ropke S, Pisinger D. An adaptive large neighborhood search heuristic
                 for the pickup and delivery problem with time windows[J].
                 Transportation Science, 2006, 40(4): 455--472.
  % ★ 若使用了 AI 工具，请在此逐条列出（示例格式，按当年官方要求调整）：
  % \bibitem{ai1} OpenAI. ChatGPT (GPT-4o) [Large language model]. 2025.
  %               https://chat.openai.com
\end{thebibliography}

% =============================================================================
% APPENDICES（计入 25 页 —— 与国赛完全相反！）
% -----------------------------------------------------------------------------
%  ★★★ 再次强调 ★★★
%   国赛：附录页数不限，且官方要求放"全部完整可运行的源代码"。
%   美赛：Appendix、Code、Table of Contents、Reference List **全部计入** 25 页。
%   因此**不要**把国赛那种长篇代码附录照搬到美赛论文里。
%   只放评阅人确实需要看的最小内容：关键推导、补充图表、极短的伪代码。
%  ★ 不要把程序文件、数据文件作为独立附件随论文一起提交
%    （官方要求只提交一个 PDF）。
% =============================================================================
\newpage
\appendix
\section{Supporting Derivations}
\label{app:derivations}

Put only what the reviewer needs; anything optional costs you pages from the 25-page
budget.

\section{Code Fragment (Optional)}
\label{app:code}

% ★ 再次提醒：这一段代码**会计入 25 页**。能删就删。
\begin{lstlisting}[language=Python,caption={Core solver loop (abridged)},label={lst:solver}]
import numpy as np
from scipy.optimize import milp, LinearConstraint, Bounds


def solve(cost, demand, capacity, fixed_cost):
    m, n = capacity.shape
    c = np.concatenate([cost.ravel(), fixed_cost.ravel()])
    lb = np.zeros(c.size)
    ub = np.concatenate([np.full(m * n, np.inf), np.ones(m)])
    A_eq = np.zeros((n, c.size))
    for j in range(n):
        A_eq[j, j::n] = 1.0
    res = milp(c=c,
               constraints=[LinearConstraint(A_eq, demand, demand)],
               integrality=np.ones(c.size),
               bounds=Bounds(lb, ub))
    return res
\end{lstlisting}

% =============================================================================
% REPORT ON USE OF AI
% -----------------------------------------------------------------------------
%  ★ 位置：**附在 25 页正文之后**。
%  ★ 页数：该部分**无页数限制，且不计入 25 页上限**。
%  ★ 透明性要求：若使用了任何 AI 工具，
%      (1) 必须在**正文中内联引用**（例如 "we used ChatGPT \cite{ai1} to ..."）；
%      (2) 必须在此处**列出所有用到的 AI 工具**及其用途；
%      (3) 必须在 References 中给出对应条目。
%  ★ 若完全没有使用 AI 工具，也要保留本节并明确声明 "No AI tools were used."
%    （按当年 Instructions 的措辞要求书写）。
%  ★ 提醒：本节的页数口径同样可能随年份调整，请以当年官方通知为准。
% =============================================================================
\newpage
\section*{Report on Use of AI}
\addcontentsline{toc}{section}{Report on Use of AI}

% 若未使用任何 AI 工具，保留下面这一段：
No generative AI tools were used in the preparation of this report.

% 若使用了 AI 工具，删除上面那段，改用下面这种逐条列出的格式：
% \begin{itemize}
%   \item \textbf{Tool:} ChatGPT (GPT-4o), OpenAI. \textbf{Use:} language polishing of
%         Section~\ref{sec:intro} and code debugging in Appendix~\ref{app:code}.
%         \textbf{Verification:} all output was reviewed and edited by the team; the
%         final models, results, and conclusions are our own.
%   \item \textbf{Tool:} DeepSeek-V3, DeepSeek. \textbf{Use:} literature search
%         suggestions. \textbf{Verification:} every suggested reference was retrieved
%         and read by the team before citation.
% \end{itemize}

\end{document}
